\centerline{\bf \Huge Домашнее задание.}
\centerline{\bf \Huge Индукция}

\begin{flushright}
        \scriptsize{Все задачи из листочка нужно решать по индукции!}
\end{flushright}

\problem{1!}
\textbf{Ханойские башни.} Есть три стержня и несколько колец разного размера. Класть можно только кольцо меньшего размера на кольцо большего размера. Можно ли переместить пирамидку с одного стрежня на другой, если в пирамидке n колец.

\problem{2!}
Докажите тождество $1^2 + 2^2 + \ldots + n^2 = \dfrac{n(n+1)(2n+1)}{6}$


\problem{3}
На плоскости проведено несколько прямых. Докажите, что части, на которые эти прямые делят плоскость, можно раскрасить в два цвета так, чтобы любые две соседние части были окрашены в разные цвета.

\problem{4}
Докажите тождество $1^3 + 2^3 + \ldots + n^3 = (1 + 2 + \ldots + n)^2$

\problem{5$^*$}
У Очира в распоряжении $2n$ цифр, среди которых есть $n$ единиц и $n$ двоек. Докажите, что он всегда может составить из них число, кратное $2^n$.